\chapter{Refinement of CRE Enforcement}
\label{chapter:refinement-of-cre-enforcement}

%%%%%%%%%%%%%%%%%%%%%%%%%%%%%%%%%%%%%%%%%%%%%%%%%%%%%%%%%%%%%%%%%%%%%%%%%%%%%%%%
%                                                                              %
%%%%%%%%%%%%%%%%%%%%%%%%%%%%%%%%%%%%%%%%%%%%%%%%%%%%%%%%%%%%%%%%%%%%%%%%%%%%%%%%
\section{Introduction}
\label{section:cre-refinement-introduction}
In this chapter we will introduce a different definition of traces (system traces)
and show that, under some conditions, it is isomorph to the (MFOTL) trace
definition introduced in chapter \ref{chapter:preliminaries}.
We will also define what \emph{Refinement} for MFOTL formulae as well as
Lex specifications.
We will show that with our definitions, if a Lex specification A is a refinement
of a Lex specification B, then A compiled to MFOTL is
an MFOTL refinement of B compiled to MFOTL.
We will further introduce a model of runtime enforcement as an event system.
This new model will be based on the new trace definition.
We will give definitions for soundness and transparency of runtime
enforcement event systems.
We will show that these definitions are equivalent to those of
chapter \ref{chapter:preliminaries} under the aforementioned isomorphism
between system and MFOTL traces.
We will adapt the definition of refinement from Sprenger et al. \cite{igloo}
slightly.
Specifically, our mediator function will not only rely on an event being reported
at a certain point in time, but also the trace (history) of previous events.

We will then go on to show that, whenever there exists a sound enforcement
event system for some MFOTL formula $\phi$, and another formula $\psi$ refines
$\phi$, then there exists a sound enforcement event system for $\psi$.
Separately we will show that there exists a sound enforcement event system
for some MFOTL formula $\phi$ if and only if there exists a sound 
enforcer, as defined in chapter \ref{chapter:preliminaries} for $\phi$.
Hublet et al. \cite{lex} have shown that an causable Lex specification
compiles to an causable MFOTL formula.
And Hublet et al. \cite{whyenf} have shown that if an MFOTL formula is 
causable, there exists a sound enforcer for it.
Putting everything together we will have shown that if a Lex specification A
refines an enforceable Lex specification B, then B is enforceable.


\section{Alternative Enforcement Model}
\label{section:updated-enforcement-model}

%%%%%%%%%%%%%%%%%%%%%%%%%%%%%%%%%%%%%%%%%%%%%%%%%%%%%%%%%%%%%%%%%%%%%%%%%%%%%%%%
\subsection{System Traces}
\label{subsection:system-traces}
We introduce a different trace definition as in Definition \ref{def:trace}
and we will show that this new definition is at least as expressive as
the old one.
This definition will help us to look at runtime enforcement
as an event system in Subsection
\ref{subsection:runtime-enforcement-event-system} that
works with traces in this new model.
By showing that the two trace definitions are isomorphic,
we will be able to show that the two
runtime enforcement models are also equivalent.
This will tie together the refinement of Lex specifications
and in turn of CRE[MFOTL] with the refinement
definition of Sprenger et al. \cite{igloo}.

\begin{definition}
    \label{def:system-sequence}
    A \emph{System Sequence} $\sigma_\sys \in \SeqSys$ 
    is a finite or infinite sequence of events
    \begin{equation*}
        \Sequence[1 \leq i \leq k]{e_i},
        k \in \Nat \cup \lbrace \infty \rbrace
    \end{equation*}
    where for every $i \in [1, k]$,
    \begin{equation*}
        e_i \in  \Databases \cup \left\{ \tick \right\}.
    \end{equation*}
\end{definition}
The special event $\tick$ represents a clock tick in the system.
As a more compact notation we use
$$e^k := \underbrace{e \cdot \ldots \cdot e}_{k \text{ times}}$$ for all $k \in \Nat$
and $e \in \Databases \cup \{tick\}$. Note that for $k=0$, $e^k = \emptytrace_\sys$.

\begin{definition}
    \label{def:num-ticks}
    The function
    \vspace{-1em}
    \begin{equation*}
        \numTicks : \SeqSys \to \Nat \cup \{ \infty \}
    \end{equation*}
    \vspace{-2em}
    \begin{align*}
        \numTicks(\sigma_\sys) =
        \begin{cases}
            % 0 &\text{if } \sigma_\sys = \emptytrace_\sys \\
            k
                &\text{if } \sigma_\sys = \tick^k
                \text{ for } k \geq 0 \\
            k + \numTicks(\sigma'_\sys)
                &\text{if } \sigma_\sys = \tick^k \cdot D \cdot \sigma'_\sys
                \text{ for } k \geq 0
                \text{ and } D \in \Databases
        \end{cases}
    \end{align*}
    gives the number of ticks in a system sequence.
\end{definition}
The $\tick$ function is useful to determine a timestamp in a
system trace, since they are only implicit and not
explicit like in MFOTL traces.

\begin{definition}
    \label{def:trim-system-traces}
    The function
    \vspace{-1em}
    \begin{equation*}
        \trim: \SeqSys \to \SeqSys
    \end{equation*}
    \vspace{-2em}
    \begin{align*}
        \trim(\sigma_\sys) :=
        \begin{cases}
            \emptytrace_\sys
                &\text{if } \sigma_\sys = \tick^k
                \text{ for } k \geq 0 \\
            \tick^k \cdot D \cdot \trim(\sigma_\sys')
                &\text{if } \sigma_\sys = \tick^k \cdot D \cdot \sigma_\sys'
                \text{ for } k \geq 0
                \text{ and } D \in \Databases
        \end{cases}
    \end{align*}
    removes all $\tick$ events from the end of a system sequence
    that are not followed by a database.
\end{definition}

For sets of system sequences $A \subseteq \SeqSys$, 
$\trim(A) := \{ \trim(\sigma) \mid \sigma \in A \}$.

\begin{definition}
    \label{def:system-trace}
    A \emph{System Trace}
    $\sigma_\sys \in \Traces_\sys \subseteq \SeqSys$ 
    is a finite or infinite system sequence where
    the number of ticks progresses with the number of databases
    (if $\numDBs(\sigma_\sys) = \infty$, then
    $\numTicks(\trim(\sigma_\sys)) = \infty $).
\end{definition}

The distinction between system sequences and system
traces is made, such that for infinite traces, it
is guaranteed that at any point in the trace, there
will eventually follow a database event.

\begin{definition}
    \label{def:num-db}
    The function
    \vspace{-1em}
    \begin{equation*}
        \numDBs : \Traces_\sys \to \Nat \cup \{ \infty \}
    \end{equation*}
    \vspace{-2em}
    \begin{align*}
        \numDBs(\sigma_\sys) =
        \begin{cases}
            0 &\text{if } \sigma_\sys = \tick^k \text{ and } k \geq 0 \\
            1 + \numDBs(\sigma'_\sys)
                &\text{if } \sigma_\sys = \tick^k \cdot D \cdot \sigma'_\sys
                \text{ for } k \geq 0
                \text{ and } D \in \Databases
        \end{cases}
    \end{align*}
    gives the number of database events in a system trace.
\end{definition}


We say a system trace $\sigma_\sys \in \Traces_\sys$
is finite iff $\numDBs(\sigma_\sys) \in \Nat$
and it is infinite iff $\numDBs(\sigma_\sys) = \infty$.
We use
$\Traces_{\sys f} \subseteq 
\left( \Databases \cup \left\{\tick \right\} \right)^*$,
$\Traces_{\sys \omega} \subseteq
\left( \Databases \cup \left\{\tick \right\} \right)^\infty$,
and
$\Traces_\sys = 
\Traces_{\sys f} \cup \Traces_{\sys \omega}$
to denote the sets of finite system traces, infinite system traces,
and all system traces respectively.
$\emptytrace_\sys$ denotes an empty system
trace.

\begin{lemma}
    \label{lem:trimmed-infinite-system-traces}
    For any infinite system trace $\sigma_\sys \in \Traces_{\sys \omega}$,
    $\sigma_\sys = \trim(\sigma_\sys).$
\end{lemma}
This follows directly from our definition for infinite system traces.
At any point within an infinite trace, there must eventually follow
at least one database, because otherwise there would not
be an infinite number of databases in the trace.
Lemma \ref{lem:trimmed-infinite-system-traces} can be rewritten as
$\Traces_{\sys \omega} = \trim(\Traces_{\sys \omega})$.

Next we define conversion functions between trimmed
system traces and MFOTL traces.
And then we will show that the conversion functions
together define an isomorphism between
MFOTL traces and trimmed system traces.
The limitation on trimmed system traces is necessary,
because MFOTL traces do not capture any timepoints at
which no database is logged.
\begin{definition}
    \vspace{-1em}
    \label{def:cconv'}
    The \emph{Finite Extended System Trace Converter Function} converts system
    traces starting at some timestamp to MFOTL traces.
    \begin{align*}
        \cConv':
        \Nat \times \Traces_\sys
        \to \Traces
    \end{align*}
    \vspace{-2em}
    \begin{align*}
        % \label{eq:cconv'}        
        \cConv'(\tau, \sigma_\sys) :=
        \begin{cases}
            \emptytrace
                % & \text{if } \trim(\sigma_\sys) = \emptytrace_\sys \\
                & \text{if } \trim(\sigma_\sys) = 0 \\
            (\tau + k, D) \cdot \cConv'(\tau+k, \sigma'_\sys)
                & \text{if } \sigma_\sys = \tick^k \cdot D \cdot \sigma'_\sys \\
                & \text{with } k \geq 0, D \in \Databases \\
        \end{cases}
    \end{align*}
\end{definition}

\begin{definition}
    \label{def:cconv}
    The \emph{System Trace Converter Function} converts system
    traces starting at timestamp $0$ to MFOTL traces.
    \vspace{-1em}
    \begin{equation*}
        \cConv: \Traces_\sys \to \Traces
    \end{equation*}
    \vspace{-2em}
    \begin{align*}
        % \label{eq:cconv}
        \cConv(\sigma_\sys) := \cConv'(0, \sigma_\sys).
    \end{align*}
\end{definition}


\begin{definition}
    \label{def:dconv'}
    The extended \emph{MFOTL Trace Converter Function}
    converts MFOTL
    traces starting at some timestamp to system traces.
    \vspace{-1em}
    \begin{align*}
        \dConv':
        \Nat \times \Traces \to \trim(\Traces_\sys)
    \end{align*}
    \vspace{-2em}
    \begin{align*}
        \label{eq:dconv'}
        \dConv'(\tau, \sigma) :=
        \begin{cases}
            \emptytrace_\sys
                &\text{if } \sigma = \emptytrace \\
            \tick^k \cdot D \cdot \dConv'(\tau+k, \sigma')
                &\text{if } \sigma = (\tau+k, D) \cdot \sigma' \\
                &\text{with } k \geq 0
        \end{cases}
    \end{align*}
\end{definition}

\begin{definition}
    \label{def:dconv}
    The \emph{MFOTL Trace Converter Function} converts MFOTL
    traces starting at timestamp $0$ to system traces.
    \vspace{-1em}
    \begin{equation*}
        \dConv: \Traces \to \trim(\Traces_\sys)
    \end{equation*}
    \vspace{-2em}
    \begin{equation*}
        \label{eq:dconv}
        \dConv(\sigma) := \dConv'(0, \sigma).
    \end{equation*}
\end{definition}

It is important to note that $\dConv'(\tau, \sigma)$
is only defined, if
$\tau \leq \fts(\sigma)$ or $\sigma = \emptytrace$.
$\dConv$ however, is defined for all $\sigma \in \Traces$,
since timestamps in any trace must increase monotonically,
and the timestamp argument passed to $\dConv'$ is $0$.
To see that the images of both $\dConv$ and $\dConv'$
for all allowed inputs is $\trim(\Traces_\sys)$ and
not $\Traces_\sys$, note that a $\tick$ is only added when
the trace is not yet empty, thus any inserted $\tick$ will
eventually be followed by some database $D$.
For any set of MFOTL traces $A \subseteq \Traces$ and any
set of system traces $B \in \Traces_\sys$ and timestamp
$\tau \in \Nat$, we write
\begin{align*}
    &\cConv(A) &&:= \{ \cConv(\sigma_\sys) \mid \sigma_\sys \in A \}, \\
    &\cConv'(\tau, A) &&:= \{ \cConv'(\tau, \sigma_\sys) \mid \sigma_\sys \in A \}, \\
    &\dConv(B) &&:= \{ \dConv(\sigma) \mid \sigma \in B \}, and \\
    &\dConv'(\tau, B) &&:= \{ \dConv'(\tau, \sigma) \mid \sigma \in B \}. \\
\end{align*}

\begin{theorem}
    \label{theo:trace-equivalence}
    The trace converter functions $\cConv$ and $\dConv$ define a
    bijective mapping between MFOTL traces $\Traces$ and
    trimmed system traces $\trim(\Traces_\sys)$
    and are each others inverses.
    \begin{equation*}
        \forall \sigma \in \Traces .~
        \sigma = \dConv(\cConv(\sigma))
    \end{equation*}
    and
    \begin{equation*}
        \forall \sigma_\sys \in \trim(\Traces_\sys) .~
        \sigma_\sys = \cConv(\dConv(\sigma))
    \end{equation*}
\end{theorem}

The intuition behind the necessity of trimming system traces to make
them fully comparable with MFOTL traces is that MFOTL traces only capture
some last timestamp at which a database was logged, while system traces
can capture a current timestamp (the total number of ticks in a system trace)
which may be greater than the last time stamp at which a database was logged
(the number of ticks up until the last database).

\begin{proof}
    \label{proof:trace-equivalence}
    The proof for Theorem \ref{theo:trace-equivalence}
    proceeds by splitting it into two properties, one
    for each direction.
    We want to show the two following equations.
    \begin{align}
        \label{eq:domain-mfotl-traces}
        &\forall \sigma \in \Traces .~
        &\sigma
        &=\dConv(\cConv(\sigma)), \\
        \label{eq:domain-system-traces}
        &\forall \sigma_\sys \in \trim(\Traces_{\sys}) .~
        &\sigma_\sys 
        &=\cConv(\dConv(\sigma_\sys)),
    \end{align}
    Since $\cConv$ and $\dConv$ are defined with
    $\cConv'$ and $\dConv'$ respectively,
    the proof will proceed by showing the following
    two, more general, properties
    for all $\tau \in \Nat$.
    \begin{align}
        \label{eq:domain-mfotl-traces'}
        &\forall \sigma \in \Traces .~
        ( \sigma = \emptytrace \vee  \tau \leq \fts(\sigma) )
        \to 
        \cConv'(\tau, \dConv'(\tau, \sigma)) = \sigma, \\
        \label{eq:domain-system-traces'}
        &\forall \sigma_\sys \in \trim(\Traces_{\sys}) .~
        \dConv'(\tau, \cConv'(\tau, \sigma_\sys)) = \sigma_\sys,
    \end{align}
    When fixing $\tau=0$ in these two equations respectively,
    the precondition in Equation \ref{eq:domain-mfotl-traces'}
    evaluates to true, and the two equations become
    equivalent to Equation \ref{eq:domain-mfotl-traces}
    and Equation \ref{eq:domain-system-traces'}.
    Thus Equation \ref{eq:domain-mfotl-traces'} implies
    Equation \ref{eq:domain-mfotl-traces} and
    Equation \ref{eq:domain-system-traces'} implies
    Equation \ref{eq:domain-system-traces}.

    For each direction we will first show that they
    hold for finite traces and then make the argument
    they also hold for infinite traces.

    \begin{enumerate}
        \item \textbf{MFOTL Traces}
            (Equations \ref{eq:domain-mfotl-traces} and
            \ref{eq:domain-mfotl-traces'})
        
        As stated above, we begin by showing Equation
        \ref{eq:domain-mfotl-traces'} for finite traces
        $\sigma \in \Traces_f$.
        In a second step we will then make the argument why it
        also holds for infinite MFOTL traces.
        We proof this by structural induction over traces $\sigma \in \Traces_f$
        for all $\tau \in \Nat$ where either $\tau \leq \fts(\sigma)$ or
        $\sigma = \emptytrace$.
    
        \textbf{Base Case:}
        $\sigma = \emptytrace$ (and some $\tau \in \Nat$)
        \begin{align*}
            \cConv'(\tau, \dConv'(\tau, \sigma))
            = \cConv'(\tau, \dConv(\tau, \emptytrace))
            \stackrel{\ref{def:dconv'}}{=}
                \cConv'(\tau, \emptytrace_\sys)
            \stackrel{\ref{def:cconv'}}{=}
                \emptytrace = \sigma
        \end{align*}
        \textbf{Induction Hypothesis:} 
        For every true suffix $\sigma'$ of a trace $\sigma \in \Traces_f$
        that is not the entire trace $\sigma$, i.e.
        $\sigma = \sigma'' \cdot \sigma'$ and $\sigma'' \neq \emptytrace$,
        and for every timestamp $\tau \in \Nat$
        which is less than the first timestamp
        ($\tau \leq \fts(\sigma')$)
        in $\sigma'$, unless $\sigma' = \emptytrace$ it holds that
        \begin{equation*}
            \sigma' = \cConv'(\tau, \dConv'(\tau, \sigma')).
        \end{equation*}
    
        \textbf{Induction Step:}
        Let $\sigma \neq \emptytrace$
        and $\tau \leq \fts(\sigma)$.
        Because $\sigma$ is not empty,
        $\sigma = \left(\tau+k, D \right) \cdot \sigma'$
        for some $k \geq 0$, some $D \in \Databases$ and
        a suffix $\sigma' \in \Traces_f$.
        \begin{align*}
            \cConv'(\tau, \dConv'(\tau, \sigma))
            &= \cConv'(\tau, \dConv'(\tau, (\tau+k, D) \cdot \sigma')) \\
            &\stackrel{\ref{def:dconv'}}{=}
                \cConv'(\tau, \tick^k \cdot D \cdot \dConv'(\tau + k, \sigma')) \\
            &\stackrel{\ref{def:cconv'}}{=}
                (\tau+k, D) \cdot \cConv'(\tau+k, \dConv'(\tau+k, \sigma')) \\
            &\stackrel{\text{IH}}{=} \left(\tau+k, D \right) \cdot \sigma'
            = \sigma
        \end{align*}

        Finally we show how this extends to infinite traces.
        Let $\sigma \in \Traces_\omega$ be an infinite MFOTL
        trace.
        And let $\sigma_{|\dots i}$ be the finite prefix with the
        first $i$ time points of $\sigma$.
        With the finite case, we have shown that
        for any $i \in \Nat$, $\sigma \in \Traces_\omega$,
        and $\tau \leq \fts(\sigma_{|\dots i})$
        (in the infinite case, a first timestamp always exists)
        $ \sigma_{..i} =
            \cConv'(\tau, \dConv'(\tau, \sigma_{|\dots i}))$.
        Then in the limit
        \begin{equation*}
            \sigma =
            \lim_{i \to \infty}
            \sigma_{|\dots i} = 
            \lim_{i \to \infty}
            \cConv'(\tau, \dConv'(\tau, \sigma_{|\dots i})) =
            \cConv'(\tau, \dConv'(\tau, \sigma)).
        \end{equation*}
    
        We have now shown \ref{eq:domain-mfotl-traces} 
        for both finite and infinite traces,
        which directly implies \ref{eq:domain-mfotl-traces'}.
        
        \item \textbf{System Traces}
            (Equations \ref{eq:domain-system-traces} and 
            \ref{eq:domain-system-traces'})
        For this direction we first prove that Equation
        \ref{eq:domain-system-traces'} holds for finite
        system traces $\sigma_\sys \in \trim(\Traces_{\sys f})$.
        And in a second step will show that it also holds
        for infinite system traces $\sigma_\sys \in \trim(\Traces_{\sys \omega})$.
        We recall that whether or not a system trace
        is finite or infinite depends on the number of
        database events in the trace.

        This first part of the proof proceeds
        by structural induction over finite and trimmed system traces
        $\sigma_\sys \in \trim(\Traces_{\sys f})$ for all $\tau \in \Nat$ .
            \begin{enumerate}
                \item
                \textbf{Base Case} $\sigma_\sys = \emptytrace_\sys$
                \begin{align*}
                    \dConv'(\tau, \cConv'(\tau, \sigma_\sys))
                    =  \dConv'(\tau, \cConv'(\tau, \emptytrace_\sys))
                    \stackrel{\ref{def:cconv'}}{=}
                        \dConv'(\tau, \emptytrace)
                    \stackrel{\ref{def:dconv'}}{=}
                        \emptytrace_\sys
                    = \sigma_\sys
                \end{align*}
                \item \textbf{Induction Hypothesis}
                For every true suffix $\sigma_\sys'$ of a trimmed system trace
                $\sigma_\sys \in \trim(\Traces_{\sys f})$
                that is not the entire trace $\sigma_\sys$, i.e.
                $\sigma_\sys = \sigma_\sys'' \cdot \sigma_\sys'$ and
                $\sigma_\sys'' \neq \emptytrace$,
                and for every timestamp $\tau \in \Nat$
                it holds that
                \begin{equation*}
                    \sigma_\sys'
                    = \dConv'(\tau, \cConv'(\tau, \sigma_\sys')).
                \end{equation*}

                \item \textbf{Induction Step}

                Let $\sigma_\sys \in \trim(\Traces_{\sys f})$,
                with $\sigma_\sys \neq \emptytrace_\sys$ .
                Then $\sigma_\sys = \tick^k \cdot D \cdot \sigma_\sys'$
                for some $k \geq 0$, $D \in \Databases$,
                and a suffix $\sigma_\sys' \in \trim(\Traces_{\sys f})$.
                \begin{align*}
                    \dConv'(\tau, \cConv'(\tau, \sigma_\sys)) &=
                    \dConv'(\tau, \cConv'(\tau, \tick^k \cdot D \cdot \sigma_\sys')) \\
                    &\stackrel{\ref{def:cconv'}}{=}
                        \dConv'(\tau, (\tau+k, D) \cdot \cConv'(\tau+k, \sigma_\sys')) \\
                    &\stackrel{\ref{def:dconv'}}{=}
                        \tick^k \cdot D \cdot
                        \dConv'(\tau+k, \cConv'(\tau+k, \sigma_\sys')) \\
                    &\stackrel{IH}{=}
                        \tick^k \cdot D \cdot \sigma_\sys'
                    = \sigma_\sys
                \end{align*}

                Finally we show how this extends to infinite system
                traces.
                Let $\sigma_\sys \in \trim(\Traces_{\sys \omega})$ be an
                infinite system trace.
                And let $\sigma_{\sys | \dots i}$ be the finite prefix
                of $\sigma_\sys$ ending with the $i$-th
                database event.
                By this definition
                $\sigma_{\sys |\dots i}$ is in $\trim(\Traces_{\sys f})$.
                With the finite case, we have shown that
                for any $i \in \Nat$, $\sigma_\sys \in \trim(\Traces_{\sys \omega})$,
                and $\tau \in \Nat$
                $ \sigma_{\sys |\dots i} =
                    \dConv'(\tau, \cConv'(\tau, \sigma_{\sys | \dots i}))$.
                Then in the limit
                \begin{equation*}
                    \sigma_\sys =
                    \lim_{i \to \infty}
                    \sigma_{\sys |\dots i} =
                    \lim_{i \to \infty}
                    \dConv'(\tau, \cConv'(\tau, \sigma_{\sys | \dots i})) =
                    \dConv'(\tau, \cConv'(\tau, \sigma_\sys)).
                \end{equation*}
                
            \end{enumerate}
        
    \end{enumerate}

\end{proof}
Intuitively $\cConv$ and $\cConv'$ count the clock ticks and store them
explicitly as integer timestamps.
% We believe that it is clear which definition is used based on the context.
% At times we may use the function $\cConv$ to make it more explicit.
For both trace definitions, a property is a subset of all possible traces.
For conversion of properties we use $\cConv$ with arguments being sets of
traces in which case $\cConv$ simply gets mapped on each element of the set.
As in Sprenger et. al \cite{igloo} and Abadi and Lamport \cite{refinement-mappings}
traces of this second form can be extended with a stuttering step
$\sskip$, where traces are equivalent if they are the same with all
$\sskip$ steps removed.
In particular this means that sets of system traces (such as trace properties)
are closed under the addition and removal of $\sskip$ events.
Thus the set of events in system traces $\Traces_\sys$
gets extended to be
$$
    \Databases \cup \left\{ \tick, \sskip \right\}.
$$
Since $\sskip$ events do not meaningfully modify a system trace,
the isomorphism from Theorem \ref{theo:trace-equivalence} works if
we remove all $\sskip$ events before doing the trace conversion.

With the isomorphism between MFOTL traces as defined in prior work and
Section \ref{chapter:preliminaries} and trimmed system traces we can
define MFOTL semantics for infinite system traces as followed.

\begin{example}
    \label{ex:system-trace}
    If we recall the MFOTL trace $\sigma$ from
    Example \ref{ex:trace}, we can now convert it into
    a system trace $\sigma_\sys := \cConv(\sigma)$.
    \begin{align*}
        \sigma_\sys = \langle
             & \{\mathsf{giveEmailConsent}(\text{"Alice"}, \text{"Account Updates"})\},
             \tick, \\ 
             & \{\mathsf{giveEmailConsent}(\text{"Bob"}, \text{"Marketing"})\},
             \tick, \\
             &  \{\mathsf{revokeEmailConsent}(\text{"Alice"}, \text{"Account Updates"}), \\
             &\quad \mathsf{giveEmailConsent}(\text{"Bob"}, \text{"Account Updates"})\}, \\
             & \tick, \tick, \tick, \tick, \tick, \tick, \tick, \tick, \emptyset \\
        \rangle \\
        % \sigma =
        % \langle
        % &(0, &&
        % \{\mathsf{giveEmailConsent}(\text{"Alice"}, \text{"Account Updates"})\}),
        % \\
        % &(1, &&
        % \{\mathsf{giveEmailConsent}(\text{"Bob"}, \text{"Marketing"})\}), \\
        % &(2,
        % &&\{\mathsf{revokeEmailConsent}(\text{"Alice"}, \text{"Account Updates"}), \\
        %      &&&\mathsf{giveEmailConsent}(\text{"Bob"}, \text{"Account Updates"})\}), \\
        % &(10, &&\emptyset)
        % \rangle
    \end{align*}
\end{example}

\begin{definition}
    \label{def:mfotl-semantics-for-system-traces}
    The semantics of MFOTL for a fixed, infinite system trace
    $\sigma_\sys \in \Traces_{\sys \omega}$, a valuation $v$,
    time-point $i$, and MFOTL formula $\phi$ is defined as
    \begin{equation*}
        v, i \vDash_{\sigma_\sys} \phi \quad \text{ iff } \quad
        v, i \vDash_{\cConv(\sigma_\sys)} \phi.
    \end{equation*}
\end{definition}

%%%%%%%%%%%%%%%%%%%%%%%%%%%%%%%%%%%%%%%%%%%%%%%%%%%%%%%%%%%%%%%%%%%%%%%%%%%%%%%%
\subsection{Runtime Enforcement as an Event System}
\label{subsection:runtime-enforcement-event-system}


In this subsection we define an alternative runtime
enforcement model than that presented in Section
\ref{subsection:proactive-runtime-enforcement}.
and by Hublet et al. \cite{whyenf}.
We want to describe runtime enforcement as an
event system, such that we can show that our
definition of CRE[MFOTL] refinement corresponds to
the notion of refinement for event systems as described
by Sprenger et al. \cite{igloo}.
Describing runtime enforcement as an event system gives
a model comparable to the edit automata approach by
Ligatti et al. \cite{edit-automata}, since both
edit automata and event systems are forms of state machines.

The event system model of runtime enforcement
works with system traces as introduced in Section
\ref{subsection:system-traces}.
At the end of this seciont, 
with the isomorphism between MFOTL traces and
trimmed system traces from Theorem
\ref{theo:trace-equivalence}, we will show that
the two runtime enforcement models are equivalent.

Our model of proactive enforcement as an event system relies on two 
update functions, one for the proactive and one for the reactive case

\begin{definition}
    \label{def:reactive-update-function}
    A \emph{Reactive Update Function}
    \begin{equation*}
       \RFun : \Traces_{\sys f} \times \Databases \to \Databases
    \end{equation*}
    modifies a reported  set of events by suppressing
    (deleting) and causing (inserting) certain events
    if needed.
    Where for all $\sigma \in \Traces_{\sys f}$
    and $D \in \Databases$
    \begin{itemize}
        \item $D \setminus \RFun(\sigma, D) \in \Databases_\Sup$
        \item $\RFun(\sigma, D) \setminus D \in \Databases_\Cau$
    \end{itemize}
\end{definition}

\begin{definition}
    \label{def:proactive-update-function}
    A \emph{Proactive Update Function}
    \begin{equation*}
       \PFun :
       \Traces_{\sys f}
       \to
       \Databases_\Cau \cup \{ \sskip \}
    \end{equation*}
    indicates whether or not some events must be caused (inserted)
    before a clock tick happens.
\end{definition}

\begin{definition}
    \label{def:enf-state-space}
    The \emph{Enforcement State Space} $\Statespace = \Traces_\sys$
    consists of all system traces (that may include $\sskip$).
\end{definition}

\begin{definition}
    \label{def:enf-initial-states}
    The single initial state is the empty trace $\emptytrace_\sys$.
    $$\Initialstates = \lbrace \emptytrace_\sys \rbrace$$
\end{definition}

\begin{definition}
    \label{def:enf-event-space}
    The \emph{Event Space}
    \begin{equation*}
        \Eventspace=\Databases \cup \lbrace \tick, \sskip \rbrace
    \end{equation*}
    consists of sets of events or the two special events $\tick$ and $\sskip$.
\end{definition}


\begin{definition}
    \label{def:enf-transition-relation}
    The \emph{Transition Relation} is a subset
    \begin{equation*}
        \Transitionrelation{}{} \subseteq
        \Traces_\sys
        \times (\Databases \cup \{\tick, \sskip \} )
        \times \Traces_\sys
    \end{equation*}
    and it is defined by the following three transition
    rules.
    \begin{align*}
        % \sigma &\Transitionrelation{}{\sskip} \sigma \cdot \sskip
        %     & \text{(Skip)} \\
        \sigma &\Transitionrelation{}{\RFun(\sigma, D)} \sigma \cdot \RFun(\sigma, D)
            & \text{(Step-R)} \\
        \sigma &\Transitionrelation{}{\PFun(\sigma) \cdot \tick}
            \sigma \cdot \PFun(\sigma) \cdot \tick
            & \text{(Step-P)}
    \end{align*}
\end{definition}
% Note that the "Skip" transition is equivalent to
% $\sigma \Transitionrelation{}{\sskip} \sigma$, since system traces
% are closed under the addition (and removal) of $\sskip$ events.
We write $\sigma \Transitionrelation{}{e} \sigma'$
as a shorthand for
$(\sigma, e, \sigma') \in \Transitionrelation{}{}$.
Further for events $e_1, e_2 \in \Eventspace$,
$\sigma \Transitionrelation{}{e_1 \cdot e_2} \sigma'$ iff
there exists $\sigma''$ such that
$\sigma \Transitionrelation{}{e_1} \sigma''$ and
$\sigma'' \Transitionrelation{}{e_2} \sigma'$.
This syntax can be extend to sequences of events (i.e. system traces)
$\rho_1, \rho_2 \in \Traces_\sys$,
$\sigma \Transitionrelation{}{\rho_1 \cdot \rho_2} \sigma'$ iff
there exists $\sigma''$ such that
$\sigma \Transitionrelation{}{\rho_1} \sigma''$ and
$\sigma'' \Transitionrelation{}{\rho_2} \sigma'$.

\begin{definition}
    \label{def:enforcement-event-system}
    An \emph{Enforcement Event System} 
    $\EventSystem = \left(\Statespace, \Eventspace, \RFun, \PFun \right)$
    is a tuple consisting of a state space $\Statespace := \Traces_\sys$,
    event space $\Eventspace = \Databases \cup \left\{ \tick, \sskip \right\}$,
    a reactive update function
    $\RFun : \Traces_{\sys f} \times \Databases \to \Databases$
    and a proactive update function
    $\PFun : \Traces_{\sys f} \to \Databases_\Cau \cup \left\{ \sskip \right\}$.
    Where the two update functions define a transition relation as 
    defined in Definition \ref{def:enf-transition-relation}.
\end{definition}

We use the same definition of event traces for an enforcement event system 
$\EventSystem$ with initial states $\Initialstates$ as for
event systems inin Sprenger et al. \cite{igloo} (Section \ref{section:event-systems}),
states $\Initialstates$.
Our described system has the property that if the only initial state is
the empty trace $\emptytrace_\sys$, the final state of the event system
and the event trace of the system are identical.

We define soundness notions for enforcement event systems that are analogous to
soundness of an enforcer in Subsection
\ref{subsection:proactive-runtime-enforcement}.
Sprenger et al. \cite{igloo} define traces as always being finite sequences of
events. We extend that definition to include both finite and infinite
event sequences.
The function $\EStraces(\EventSystem, \Initialstates)$ then denotes both all
finite and infinite event sequences of the event system $\EventSystem$
starting in any initial state in $\Initialstates$.
$\EStraces_f(\EventSystem, \Initialstates)$ denotes only the finite traces
and $\EStraces_\omega(\EventSystem, \Initialstates)$ only the
infinite traces.
We also extend the definition of trace properties to include both
finite and infinite sequences.
We use the following notations if we specifically talk
about trace properties that are strictly finite or infinite respectively:
$\EventSystem, \Initialstates \vDash_f P$ iff
$\EStraces_f(\EventSystem, \Initialstates) \subseteq P$
and
$\EventSystem, \Initialstates \vDash_\omega P$ iff
$\EStraces_\omega(\EventSystem, \Initialstates) \subseteq P$.
Like Sprenger et. al \cite{igloo} we will omit the set braces
for initial states, if the set is a singleton.
Note that system traces are only considered infinite if they
are infinite with all skips removed and trimmed (ticks not followed
by a database removed).

Next we will show that any enforcer as in Hublet et al. \cite{whyenf} has a
corresponding enforcement event system.
To simplify this we define a helper function, that can compute the
state of an enforcer after having enforced some trace as well as potentially
some "leftover" ticks.

\begin{definition}
    \label{def:enforcer-state-function}
    The \emph{Enforcer State Function}
    \begin{equation*}
        \state : \EnfStates \times \Traces_\sys \to \EnfStates
    \end{equation*}
    gives the state starting in some initial state, after having enforced
    the corresponding MFOTL trace of a system trace as well as the ticks
    that may or may not be "dangling" at the end of the system trace.
    This function is defined by essentially simulating the $\run$ function of
    Algorithm \ref{alg:loop}.
    The simulation uses a helper function
    \begin{equation*}
        \state' : \EnfStates \times \Traces_\sys \times \Traces_\sys \to \EnfStates
    \end{equation*}
    $\state'(s, \sigma_{\sys 1}, \sigma_{\sys 2}) := $
    \begin{align*}
        \begin{cases}
            s
                &\text{if } \sigma_{\sys 2} = \emptytrace_\sys \\
            \state'(s', \sigma_{\sys 1} \cdot D_\Cau \cdot \tick, \sigma_{\sys 2}')
                &\text{if } \sigma_{\sys 2} = \tick \cdot \sigma_{\sys 2}' \\
                &\text{and } (\PCom(D_\Cau), s') \\
                &\quad = \EnfFun(\cConv(\sigma_{\sys 1}), s, \numTicks(\sigma_{\sys 1}), \tick) \\
            \state'(s', \sigma_{\sys 1} \cdot \tick, \sigma_{\sys 2}')
                &\text{if } \sigma_{\sys 2} = \tick \cdot \sigma_{\sys 2}' \\
                &\text{and } (\NCom, s') \\
                &\quad = \EnfFun(\cConv(\sigma_{\sys 1}), s, \numTicks(\sigma_{\sys 1}), \tick) \\
            \state'(s', \sigma_{\sys 1} \cdot (D \setminus D_\Sup \cup D_\Cau), \sigma_{\sys 2}')
                &\text{if } \sigma_{\sys 2} = D \cdot \sigma_{\sys 2}' \\
                &\text{and } (\RCom(D_\Cau, D_\Sup), s') \\
                &\quad = \EnfFun(\cConv(\sigma_{\sys 1}), s, \numTicks(\sigma_{\sys 1}), D)
        \end{cases}
    \end{align*}
    Then the $\state$ function is defined as followed for any enforcer state
    $s \in \EnfStates$ and system trace~$\sigma_\sys \in \Traces_\sys$
    \begin{equation*}
        \state(s, \sigma_\sys) := \state'(s, \emptytrace_\sys, \sigma_\sys).
    \end{equation*}
\end{definition}

\begin{definition}
    \label{def:induced-event-system}
    The \emph{induced enforcement event system} of an enforcer
    $\Enforcer = (\EnfStates, s_0, \EnfFun)$ is defined as
    $ \EventSystem(\Enforcer) :=
    (\Traces_\sys, \Databases \cup \{\tick, \sskip \}, \RFun, \PFun), $
    with  $\RFun$ and $\PFun$ defined as followed
    \begin{align*}
        \RFun(\sigma_\sys, D)
        &:= D \setminus D_\Sup \cup D_\Cau \\
        &\qquad \text{with }
        \EnfFun(\cConv(\sigma_\sys), \state(s_0, \sigma_\sys), \numTicks(\sigma_\sys), D) =
        (\RCom(D_\Cau, D_\Sup), s) % TODO: how to involve underlying system state??
        \\
        \PFun(\sigma_\sys) &:=
        \begin{cases}
            \sskip &\text{if }
                \EnfFun(\cConv(\sigma_\sys), \state(s_0, \sigma_\sys), \numTicks(\sigma_\sys), D) =
                (\NCom, s) \\
            D_\Cau &\text{if }
                \EnfFun(\cConv(\sigma_\sys), \state(s_0, \sigma_\sys), \numTicks(\sigma_\sys), D) =
                (\PCom(D_\Cau), s)
        \end{cases}
    \end{align*}
\end{definition}

\begin{definition}
    \label{def:enf-event-system-soundness-property}
    An enforcement event system
    $\EventSystem =
    \left(\Traces_\sys, \Databases \cup \{ \tick, \sskip \}, \RFun, \PFun \right)$
    is \emph{sound} with respect to some
    property $P \subseteq \Traces_{\sys \omega}$ iff
    \begin{equation*}
        \EventSystem, \emptytrace_\sys \vDash P.
    \end{equation*}    
\end{definition}
Equivalently an enforcement event system $\EventSystem$ and
the initial state $\emptytrace_\sys$
are sound for a property $P \subseteq \Traces_{\sys \omega}$
iff $\EStraces_\omega(\EventSystem, \emptytrace_\sys) \subseteq P$.
This soundness definition difers from
Sprenger et al. \cite{igloo} in that
it is over infinite traces, where Sprenger et al.
define soundness over finite traces.
We however want to use infinite traces, because
the enforcement model of Hublet et al. \cite{whyenf}
defines soundenss with infinite traces.
This deviation from Sprenger et al. \cite{igloo}
is fine to make, because any set
of infinite traces could be converted into an
(infinitely large) set of all finite prefixes of
the inifinite traces.

\begin{definition}
    \label{def:enf-event-system-soundness-formula}
    An enforcement event system $\EventSystem$ is sound for a formula $\phi$
    iff it is sound for the property $\Language(\phi)$.
\end{definition}

\begin{theorem}
    \label{theo:induced-enforcement-event-system-soundness}
    If an enforcer $\Enforcer = (\EnfStates, s_0, \EnfFun)$ is sound for a
    property $P \subseteq \Traces_\omega$, then the induced enforcement event
    system $\EventSystem(\Enforcer)$ is sound for the property
    $\cConv(P) \subseteq \Traces_{\sys \omega}$.
\end{theorem}

Before we can prove Theorem
\ref{theo:induced-enforcement-event-system-soundness},
we need the following two lemmata.
The proofs for which can be found in Appendix
\ref{appendix:section:proof-splitting-lemma} and
\ref{appendix:section:proof-enf-es-trace-correspondance}.
\begin{lemma}
    \label{lem:cconv'-splitting}
    For all $\sigma_\sys$, $\tau \in \Nat$ and $D \in \Databases$,
    \begin{equation*}
        \cConv'(\tau, \sigma_\sys \cdot D)
        = \cConv'(\tau, \sigma_\sys) \cdot
        \cConv(\tau + \numTicks(\sigma_\sys), D)
    \end{equation*}
\end{lemma}

\begin{lemma}
    \label{lem:enf-es-trace-correspondance}
    Let $\Enforcer = (\EnfStates, s_0, \mu)$ be an enforcer.
    For any system trace
    $\sigma_\sys \in \EStraces(\EventSystem(\Enforcer), \emptytrace_\sys)$
    there exists a trace $\sigma \in \Traces$ such that
    $\cConv(\sigma_\sys) = \Enforcer(\sigma)$.
\end{lemma}

\begin{proof}
    \label{proof:theorem-induced-enforcement-event-system-soundness}
    We will show that for all
    $\sigma_\sys \in \EStraces_f(\EventSystem(\Enforcer), \emptytrace_\sys)$
    there exists a trace $\sigma_\sys' \in \cConv(P)$
    of which $\sigma_\sys$ is a prefix
    (i.e. $\sigma_\sys = \sigma'_{\sys \mid \ldots \leq |\sigma_\sys|}$).

    Since by definition
    $\EStraces_f(\EventSystem(\Enforcer), \emptytrace_\sys)$
    contains all finite prefixes of
    $\EStraces_\omega(\EventSystem(\Enforcer), \emptytrace_\sys)$,
    if
    $\EStraces_\omega(\EventSystem(\Enforcer), \emptytrace_\sys)
    \not\subseteq \cConv(P)$,
    there would exist some
    contains all finite prefixes of
    $\sigma_\sys \in \EStraces_f(\EventSystem(\Enforcer), \emptytrace_\sys)$
    which is not a prefix for any trace $\sigma_\sys' \in \cConv(P)$.
    Thus if we show that every finite prefix 
    $\sigma_\sys \in \EStraces_f(\EventSystem(\Enforcer), \emptytrace_\sys)$
    is also a finite prefix for at least one trace  $\sigma_\sys' \in \cConv(P)$,
    we will have shown that
    $\sigma_\sys \in \EStraces_\omega(\EventSystem(\Enforcer), \emptytrace_\sys)
    \subseteq \sigma_\sys' \in \cConv(P)$.

    This property is a consequence of Lemma
    \ref{lem:enf-es-trace-correspondance}.
    In Lemma
    \ref{lem:enf-es-trace-correspondance}
    we have shown that any enforcement event system trace
    corresponds to a trace produced by the enforcer.
    All those traces "produced" by the enforcer fall into $P$,
    by the fact that the enforcer is sound for $P$.
    and since we have shown that for all system traces
    $\sigma_\sys \in \EStraces_f(\EventSystem(\Enforcer), \emptytrace_\sys)$,
    there exists a trace $\sigma \in \Traces$
    such that $\cConv(\sigma_\sys) = \Enforcer(\sigma)$, we have
    also shown that $\cConv(\sigma_\sys) \in P$ or equivalently
    $\cConv(\EStraces(\EventSystem(\Enforcer), \emptytrace_\sys)) \subseteq P$.
    This does also imply 
    $\EStraces(\EventSystem(\Enforcer), \emptytrace_\sys) \subseteq \cConv(P)$.
\end{proof}

%%%%%%%%%%%%%%%%%%%%%%%%%%%%%%%%%%%%%%%%%%%%%%%%%%%%%%%%%%%%%%%%%%%%%%%%%%%%%%%%
%                                                                              %
%%%%%%%%%%%%%%%%%%%%%%%%%%%%%%%%%%%%%%%%%%%%%%%%%%%%%%%%%%%%%%%%%%%%%%%%%%%%%%%%
\section{Refinement}
\label{section:refinement}

%%%%%%%%%%%%%%%%%%%%%%%%%%%%%%%%%%%%%%%%%%%%%%%%%%%%%%%%%%%%%%%%%%%%%%%%%%%%%%%%
\subsection{Enforcement Event System Refinement}

For Enforcement Event Systems we will now define refinement
via a mediator function and a simulation relation as in Sprenger et. al \cite{igloo}.

\begin{definition}
    \label{def:enforcement-event-system-refinement}
    An enforcement event system
    $$\EventSystem_C =
    ( \Traces_{\sys C}, \Databases_C \cup \{\tick, \sskip \}, \RFun_C, \PFun_C )$$
    for some signature 
    $\Sigma_C = \left( \Domain, \Events_C, \Functions, \iota_C \right)$
    with initial states $\Initialstates_C \subseteq \Traces_C$
    refines an enforcement event system
    $\EventSystem_A = \left( \Traces_{\sys A}, \Databases_A \cup \{\tick, \sskip \}, \RFun_A, \PFun_A \right)$
    for some signature 
    $\Sigma_A = \left( \Domain, \Events_A, \Functions, \iota_A \right)$
    with initial states $\Initialstates_A \subseteq \Statespace_A$
    via a mediator function
    \begin{equation*}
        \MFun : \Traces_C \times (\Databases_C \cup \{ \tick, \sskip \})
        \to (\Databases_A \cup \{ \tick, \sskip \})
    \end{equation*}
    iff
    there is a simulation relation
    \begin{equation*}
        \SimRel \subseteq \Traces_C \times \Traces_A
    \end{equation*}
    such that
    \begin{enumerate}
        \item
            for all $\sigma_C \in \Initialstates_C$ there exists some
            $\sigma_A \in \Initialstates_A$ such that
            $(\sigma_C, \sigma_A) \in \SimRel$, and
        \item
            for all $\sigma_A \in \Statespace_A$,
            $\sigma_C, \sigma_C' \in \Statespace_C$, and
            $e_C \in \Eventspace_C$
            such that $(\sigma_C, \sigma_A) \in \SimRel$ and
            $\sigma_C \Transitionrelation{}{e_C} \sigma_C'$
            there exists $\sigma_A'$
            such that $(\sigma_C', \sigma_A') \in \SimRel$ and
            $\sigma_A \Transitionrelation{}{\MFun(\sigma_C, e_C)} \sigma_A'$.
    \end{enumerate}
    
\end{definition}
It is important to highlight, that our definition of a mediator function
differs from that in Sprenger et al. \cite{igloo} in that we include a second
argument for the trace up to the current point in time.
The reason for this is how we define MFOTL and Lex refinement in the following
sections.
Without going into the details, on a high level MFOTL refinements can define
let bindings for abstract events, thus "concretizing" them.
And let bindings allow for arbitrary MFOTL formulae on their right hand side which
can of course depend on the past.
That is why we need the history of events as an additional argument for our mediator
function.

\begin{definition}
    \label{def:mediator-function-trace-extension}
    A mediator function
    $$\MFun : \Traces_C \times (\Databases_C \cup \{ \tick, \sskip \})
    \to (\Databases_A \cup \{ \tick, \sskip \})$$
    can be extended to a trace transducer function
    $$\vec{\MFun} : \Traces_C \to \Traces_A$$
    as followed
    \begin{align*}
        \vec{\MFun}(\sigma_\sys) :=
        \begin{cases}
            \emptytrace_\sys &\text{if } \sigma_\sys = \emptytrace_\sys \\
            \vec{\MFun}(\sigma_\sys') \cdot \MFun(\vec{\MFun}(\sigma_\sys'), e)
                &\text{if } \sigma_\sys = \sigma_\sys' \cdot e \\
                &\text{for some } e \in \Databases_C \cup \{\tick, \sskip \}
        \end{cases}
    \end{align*}
\end{definition}

The following theorem is analogous to
Theorem \ref{theo:refinement-traces}
\cite[Theorem 2.1]{igloo}.
\begin{theorem}
    \label{theo:refinement-traces-enf-event-system}
    $(\EventSystem_C, \Initialstates_C) \refineL_\MFun (\EventSystem_A, \Initialstates_A)$
    implies
    $\vec{\MFun}(\EStraces(\EventSystem_C, \Initialstates_C)) \subseteq
    \EStraces(\EventSystem_A, \Initialstates_A)$.
\end{theorem}
The proof for Theorem \ref{theo:refinement-traces-enf-event-system} is
analogous to the proof for Theorem \ref{theo:refinement-traces}.

The following lemma is analogous to
Lemma \ref{lem:refinement-property-fullfillment}
\cite[Lemma 2.2]{igloo}.
\begin{lemma}
    \label{lem:refinement-property-fullfillment-enf-event-system}
    Suppose 
    $\EventSystem_A, \Initialstates_A \vDash P$
    and
    $\vec{\MFun}(\EStraces(\EventSystem_C, \Initialstates_C)) \subseteq
    \EStraces(\EventSystem_A, \Initialstates_A)$
    then
    $\EventSystem_C, \Initialstates_C \vDash \vec{\MFun}^{-1}(P)$
\end{lemma}
The proof for Lemma \ref{lem:refinement-property-fullfillment-enf-event-system}
is analogous to the proof for Lemma \ref{lem:refinement-property-fullfillment}.


%%%%%%%%%%%%%%%%%%%%%%%%%%%%%%%%%%%%%%%%%%%%%%%%%%%%%%%%%%%%%%%%%%%%%%%%%%%%%%%%
\subsection{CRE[MFOTL] Refinement}
\label{subsection:refinement-of-cre-specifications}
In this subsection we recall refinement of Lex specifications as it is described in
Hublet et. al \cite{lex}. We then go ahead and proof that whenever a CRE
specification refines another CRE specification, their compiled MFOTL
formulas are likewise refinements of each other.

\begin{figure}
    \begin{gather*}
        \prftree[r]
            %%begin novalidate
            {$\Events^+_\mathsf{mon}$}
            %%end novalidate
            {\mathsf{let}_d \not\in \mathsf{dom}~\Gamma}
            {\Gamma \vdash e(\bar{x}) : \mono{e}^+}
        \quad
        \prftree[r]
            %%begin novalidate
            {$\Events^\mathsf{let}_\mathsf{mon}$}
            %%end novalidate
            {\mathsf{let}_d \in \mathsf{dom}~\Gamma}
            {\Gamma(\mathsf{let}_d) = \phi} 
            {\Gamma \setminus \mathsf{let}_d \vdash \phi : \mono{e}^p}
            {\Gamma \vdash d(\bar{x}) : \mono{e}^p}
        \\[5pt]
        \prftree[r]
            %%begin novalidate
            {$\mathsf{occ}^p_\mathsf{mon}$}
            %%end novalidate
            {e \not \in \events{\phi}}
            {\Gamma \vdash \phi : \mono{e}^p}
        \quad
        \prftree[r]
            %%begin novalidate
            {$\neg^p_\mathsf{mon}$}
            %%end novalidate
            {\Gamma \vdash \phi : \mono{e}^{\neg p}}
            {\Gamma \vdash \neg \phi : \mono{e}^p}
        \\[5pt]
        \prftree[r]
            %%begin novalidate
            {$\wedge^p_\mathsf{mon}$}
            %%end novalidate
            {\Gamma \vdash \phi : \mono{e}^p}
            {\Gamma \vdash \psi : \mono{e}^p}
            {\Gamma \vdash \phi \wedge \psi : \mono{e}^p}
        \quad
        \prftree[r]
            %%begin novalidate
            {$\exists^p_\mathsf{mon}$}
            %%end novalidate
            {\Gamma \vdash \phi : \mono{e}^p}
            {\Gamma \vdash \exists x.~ \phi : \mono{e}^p}
        \\[5pt]
        \prftree[r]
            %%begin novalidate
            {$\Next^p_\mathsf{mon}$}
            %%end novalidate
            {\Gamma \vdash \phi : \mono{e}^p}
            {\Gamma \vdash \Next \phi : \mono{e}^p}
        \quad
        \prftree[r]
            %%begin novalidate
            {$\Prev^p_\mathsf{mon}$}
            %%end novalidate
            {\Gamma \vdash \phi : \mono{e}^p}
            {\Gamma \vdash \Prev \phi : \mono{e}^p}
        \\[5pt]
        \prftree[r]
            %%begin novalidate
            {$\Since^p_\mathsf{mon}$}
            %%end novalidate
            {\Gamma \vdash \phi : \mono{e}^p}
            {\Gamma \vdash \psi : \mono{e}^p}
            {\Gamma \vdash \phi \Since \psi : \mono{e}^p}
        \quad
        \prftree[r]
            %%begin novalidate
            {$\UUntil^p_\mathsf{mon}$}
            %%end novalidate
            {\Gamma \vdash \phi : \mono{e}^p}
            {\Gamma \vdash \psi : \mono{e}^p}
            {\Gamma \vdash \phi \UUntil \psi : \mono{e}^p}
        \\[5pt]
        \prftree[r]
            %%begin novalidate
            {$\mathsf{let}^p_\mathsf{mon}$}
            %%end novalidate
            {\Gamma, \mathsf{let}_d : \phi \vdash \psi : \mono{e}^p}
            {\Gamma \vdash \LLet{d(\bar{x})}{\phi}{\psi} : \mono{e}^p}
    \end{gather*}
    \vspace{-2em}
    \caption{MFOTL (Anti) Monotonicity Rules \cite[Appendix A]{lex}}
    \label{fig:mfotl-monotonicity}
\end{figure}

\begin{figure}
    \begin{gather*}
    \prftree[r] {$\textsc{con}^q_\mathsf{mon}$}
        {
            \prfStackPremises
            {\exists r : \phi \Rightarrow p \in C}
            {\forall r : \phi \Rightarrow p \in C.~
                \Gamma \vdash^\Phi \phi : \mono{e} ^q~
                \textbf{and}~\Gamma \vdash^\Phi \exc(r) : \mono{e}^{\neg q}}
        }
        {\Gamma \vdash^\Phi p(\overline{x}) : \mono{e}^q} \\[5pt]
    \prftree[r] {$\textsc{exc}^p_\mathsf{mon}$}
        {\forall r : \phi \Rightarrow \exc(s) \in E.~
        \Gamma \vdash^\Phi \phi : \mono{e}^p~
        \textbf{and}~\Gamma \vdash^\Phi \exc(r) : \mono{e}^{\neg p}}
        {\Gamma \vdash^\Phi \exc(s)(\overline{x}) : \mono{e}^p} \\[5pt]
    \prftree[r] {$\textsc{reg}^p_\mathsf{mon}$}
        {r: \phi \Rightarrow \psi \in R}
        {\Gamma \vdash \phi \to \psi : \mono{e}^p}
        {\Gamma \vdash \exc(r) : \mono{e}^p}
        {\Gamma \vdash^\Phi \comply(r) : \mono{e}^p} \\[5pt]
    \prftree {r: \phi \Rightarrow \psi \in R}
        {\Gamma \vdash (C,R,E) : \mono{\exc(r)}^+} \quad
    \prftree {r: \phi \Rightarrow p \in C}
        {\Gamma \vdash (C,R,E) : \mono{p}^{\neg q}}
        {\Gamma \vdash (C,R,E) : \mono{\exc(r)}^q} \\[5pt]
    \prftree {r: \phi \Rightarrow \exc(s) \in E}
        {\Gamma \vdash (C,R,E) : \mono{\exc(s)}^{\neg p}}
        {\Gamma \vdash (C,R,E) : \mono{\exc(r)}^p} \\[5pt]
    \prftree {\forall r: \phi \Rightarrow \psi \in R.~
        \Gamma \vdash^{(C,R,E)} \comply(r) : \mono{e}^p }
        {\Gamma \vdash (C,R,E) : \mono{e}^p}
    \end{gather*}
    \vspace{-2em}
    \caption{CRE (Anti) Monotonicity Rules \cite[Appendix A]{lex}}
    \label{fig:cre-monotonicity-rules}
\end{figure}

\begin{figure}[t]
    \begin{gather*}
        \phi \Updownarrow^+ \psi := \phi \Longrightarrow \psi\qquad
        \phi \Updownarrow^- \psi := \phi \Longleftarrow \psi \\[5pt]
        \prftree[r]
            {\textsc{refl}}
            {S,\Gamma \leadsto_{\mathrm{id}} S,\Gamma}\quad
        \prftree[r]
            {\textsc{trans}}
            {S_1,\Gamma_1 \leadsto_{\pi_1} S_2,\Gamma_2}
            {S_2,\Gamma_2 \leadsto_{\pi_2} S_3,\Gamma_3}
            {S_1,\Gamma_1 \leadsto_{\pi_1 \circ \pi_2} S_3,\Gamma_3}\\[5pt]
        \prftree[r]
            {\textsc{ref}}
            {
              \prfStackPremises
              {C' = C \cup \{r : \phi \Rightarrow e(x_1,\dots,x_k)\} \quad
                  \mathsf{events}(\phi) \setminus \Events=\{e_1,\dots,e_\ell\}}
              {\Events' = \Events \setminus \{e\} \cup \{e_1,\dots,e_\ell\} \quad
                  \iota' = \iota[i.~e_i \mapsto k_i] \setminus \{e\}}
              {\kappa' = \kappa[i.~e_i \mapsto \tau_i] \setminus \{e\} \quad
                  \Gamma' = \Gamma[i.~e_i \mapsto \tau'_i] \setminus \{e\} }
              {\forall i.~\tau'_i \sqsubseteq \tau_i \quad \Gamma' \vdash \phi : \Gamma(e)
                  \quad
                  \nexists r', \phi', \overline{x}'.~r': \phi' \Rightarrow e(\overline{x}') \in C}
              {\pi = \textsc{\textsf{map}}_{e(x_1,\dots,x_k),\phi,\{e_1,\dots,e_\ell\}}}
            }
            {((\Domain,\Events,\Functions,\iota),(C,R,E),\kappa),\Gamma
            \leadsto_\pi
            ((\Domain,\Events',\Functions,\iota'),(C',R,E),\kappa'),\Gamma'}\\[5pt]
        \prftree[r]
            {$\textsc{imp}_R$}
            {
                \prfStackPremises
                {R' = R \cup \{r : \phi' \Rightarrow \psi' \}
                    \setminus \{r : \phi \Rightarrow \psi\}}
                {\phi' \to \psi' \Longrightarrow \phi \to \psi \quad
                \Gamma \vdash^{(C,R',E)} \comply(r) : \Cau}
            }    
            {((\Domain,\Events,\Functions,\iota),(C,R,E),\kappa),\Gamma
            \leadsto_{\mathrm{id}}
            ((\Domain,\Events,\Functions,\iota),(C,R',E),\kappa),\Gamma}\\[5pt]
        \prftree[r]
            {$\textsc{imp}_C$}
            {
              \prfStackPremises
              {C' = C \cup \{ r: \phi' \Rightarrow p \}
                  \setminus \{ r : \phi \Rightarrow p \} \quad
                R' = R \cup \{ s : \psi \Rightarrow \chi \}}
              {q \in \{+,-\} \quad
                  \Always (\psi \to \chi) \Longrightarrow  \phi' \Updownarrow^q \phi}
              {(C \setminus \{ r:\phi\Rightarrow p\},R',E) : \mono{p}^q \quad
                \Gamma \vdash^{(C',R,E)} \comply(s) : \Cau}
              {\forall \tau \in \{\Cau,\Sup\}.~\Gamma
                  \vdash^{(C,R,E)} p : \tau \Longrightarrow
                  \Gamma \vdash^{(C',R,E)} p : \tau}
            }    
            {((\Domain,\Events,\Functions,\iota),(C,R,E),\kappa),\Gamma
                \leadsto_{\mathrm{id}}
                ((\Domain,\Events,\Functions,\iota),(C',R',E),\kappa),\Gamma}\\[5pt]
        \prftree[r]
            {$\textsc{imp}_E$}
            {
                \prfStackPremises
                {E' = E \cup \{ r : \phi' \Rightarrow \exc(s) \}
                    \setminus \{ r : \phi \Rightarrow \exc(s) \} \quad
                    R' = R \cup \{ s : \psi \Rightarrow \chi\}}
                {q \in \{+,-\} \quad \Always (\psi \to \chi) \Longrightarrow
                    \phi' \Updownarrow^q \phi}
                {(C,R',E \setminus \{ r:\phi\Rightarrow \exc(s)\}) : \mono{\exc(s)}^q
                    \quad
                    \Gamma \vdash^{(C,R',E')} \comply(s) : \Cau }
                {\forall \tau \in \{\Cau,\Sup\}.~\Gamma
                    \vdash^{(C,R,E)} \exc(s) : \tau \Longrightarrow
                    \Gamma \vdash^{(C,R,E')} \exc(s) : \tau}
            }
            {((\Domain,\Events,\Functions,\iota),(C,R,E),\kappa),\Gamma
                \leadsto_{\mathrm{id}}
                ((\Domain,\Events,\Functions,\iota),(C,R',E'),\kappa),\Gamma}
    \end{gather*}
    \vspace{-2em}
    \caption{CRE[MFOTL] refinement rules \label{fig:cre-refinement-rules} \cite[Figure 7]{lex}}
\end{figure}
\begin{figure}[t]
    \begin{center}
        \begin{algorithmic}[1]
            \Function{\textsf{map}$_{e(x_1,\dots,x_k),\phi,E}$}{$\sigma$}
            \State $\langle(\tau_i, D_i)_{i \in \Nat}\rangle \gets \sigma$
            \For{$i \in \Nat$}
            \State $D'_i \gets D_i \cup \{ e(v(x_1),\dots,v(x_k)) \mid
                v,i \vDash_\sigma \phi  \} $\label{alg:map:add}
            \State $D'_i \gets D'_i \setminus \{ e(\overline{d}) \mid
                (d_1,\dots,d_{\iota(e)}) \in \Domain^{\iota(e)}, e \in E \}$
                \label{alg:map:remove}
            \EndFor
            \State \Return $\langle(\tau_i, D'_i)_{i \in \Nat}\rangle$
            \EndFunction
        \end{algorithmic}
    \end{center}
    \vspace{-1.2em}
    \caption{\textsc{\textsf{map}} function used in Figure~\ref{fig:cre-refinement-rules}
    \label{fig:cre-refinement-map-function}}
\end{figure}

Hublet et al. \cite{lex} use the rules in Figure \ref{fig:cre-refinement-rules}
to define refinement for CRE[MFOTL] specifications in the following manner.
\begin{definition}
    \label{def:cre-refinement}
    Refinement of CRE specifications is defined by
    the rules in Figure \ref{fig:cre-refinement-rules}.
    A specification $S_C$ refines a specification $S_A$, written as
    $S_C \refineL_\pi S_A$
    under a refinement mapping $\pi$
    iff there exist $\Gamma_C$ and $\Gamma_A$ such that
    $S_C, \Gamma_C \leadsto_\pi S_A, \Gamma_A$.
\end{definition}
The refinement rules in Figure \ref{fig:cre-refinement-rules} depend on the concept of
monotonic and anti-monotonic predicates.
Figure \ref{fig:mfotl-monotonicity} shows the monotonicity rules for
MFOTL predicates, in which $p,q \in \{+,-\}$, with $\neg + := -$ and $\neg - := +$.
A predicate $P$ appears monotonically (or anti-monotonically) in a
formula $\phi$, if the addition of an event $P(\bar{x})$ cannot make
an already compliant (or non-compliant) trace non-compliant (or compliant).
For example in the formula $\phi := \Always (\neg p(1))$ the
predicate $P$ does not appear monotonically,
because the addition of an event at any time $P(1)$ would make a trace that
was previously compliant with $\phi$, non-compliant with $\phi$,
on the other hand $P$ appears anti-monotonically in $\phi$.
It is important to note that a predicate can appear monotonically, anti-monotonically,
both monotonically and anti-monotonically, or neither monotonically nor
anti-monotonically in a formula.

Figure \ref{fig:cre-refinement-map-function} shows an
algorithm that converts a concrete MFOTL trace
to an abstract one, according to a formula $\phi$
that defines an event $e(x_1,\dots,x_k)$
in terms of the concrete events $e_1,\dots,e_\ell$.
This function does for MFOTL traces, what the
mediator function $\pi$ from Definition
\ref{def:cre-refinement} and its extended form
$\vec{pi}$ from Definition
\ref{def:mediator-function-trace-extension} do for system traces.

\begin{definition}
    \label{def:mfotl-map-function-to-mediator-function}
    Let $\textsc{\textsf{map}}$ be a function as
    defined in Figure \ref{fig:cre-refinement-map-function}.
    Then the mediator function $\cConv(\textsc{\textsf{map}})$
    is defined as follows:
    \begin{align*}
        \cConv(\textsc{\textsf{map}})(\sigma_\sys, \tick) &:= \tick \\
        \cConv(\textsc{\textsf{map}})(\sigma_\sys, \sskip) &:= \sskip \\
        \cConv(\textsc{\textsf{map}})(\sigma_\sys, D_C) &:= D_A
    \end{align*}
    where $D_A$ is defined as the database $D'$
    being returned by the $\textsc{\textsf{map}}$ function
    for the input trace $\cConv(\sigma_\sys \cdot D)$ as
    the database for the final time point of the returned
    trace.
\end{definition}

\begin{lemma}
    \label{lem:cre-refinement-mediator-function}
    For any MFOTL trace $\sigma \in \Traces$, and
    map function $\textsc{\textsf{map}}$
    \begin{equation*}
        \overrightarrow{\cConv(\textsc{\textsf{map}})}(\cConv(\sigma))
        = \cConv(\textsc{\textsf{map}}(\sigma))
    \end{equation*}
\end{lemma}

The following theorem states that any valid refinement
of an enforceable CRE[MFOTL] specification is itself
enforceable.
\begin{theorem}
    \label{theo:cre-refinement-enforceability}
    If $S_C \refineL_\pi S_A$ under a refinement mapping
    and $S_A$ is enforceable, then $S_C$ is enforceable.
\end{theorem}

\begin{theorem}
    \label{theo:cre-refinement-trace-containment}
    If $S_C \refineL_\pi S_A$,
    $S_A$ is enforceable,
    and $\sigma \vDash S_C$,
    then $\pi(\sigma) \vDash S_A$.
\end{theorem}

Hublet et al. \cite{lex} have proved Theorem
\ref{theo:cre-refinement-enforceability}
\cite[Theorem 4]{lex} and Theorem
\ref{theo:cre-refinement-trace-containment}
\cite[Theorem 3]{lex}.

\begin{theorem}
    \label{theo:cre-refinement-induced-event-system-refinement}
    If $S_C \refineL_\pi S_A$ and $S_A$ is enforceable
    then there exist two enforcement event systems
    $\EventSystem_C$ and $\EventSystem_A$ such that
    they are sound for $S_C$ and $S_A$ respectively,
    and refine each other via the mediator function
    $\cConv(\pi)$, i.e.
    \begin{equation*}
        \EventSystem(\Enforcer_C) \refineL_{\cConv(\pi)} \EventSystem(\Enforcer_A).
    \end{equation*}
\end{theorem}

% Due to time constraints the proof for Theorem
% \ref{theo:cre-refinement-induced-event-system-refinement}
% is only outlined here, and some parts remain to be worked out
% in full detail

\begin{proof}
    for Theorem \ref{theo:cre-refinement-induced-event-system-refinement}
    Let $S_C \refineL_\pi S_A$,
    and $S_A$ be enforceable.

    By the enforceability of $S_A$ we know that
    there exists an enforcer
    $\Enforcer_A = (\EnfStates_A, s_0, \EnfFun_A)$
    which is sound for $S_A$.
    By Theorem \ref{theo:cre-refinement-enforceability}
    we also know that there exists an enforcer
    $\Enforcer_C = (\EnfStates_C, s_0, \EnfFun_C)$
    which is sound for $S_C$.

    Since $S_C \refineL_\pi S_A$ there exist two
    contexts $\Gamma_C$ and $\Gamma_A$
    such that $S_A, \Gamma_A \leadsto_\pi S_C, \Gamma_C$.
    This proof proceeds by induction over the different
    CRE[MFOTL] refinement rules for $\leadsto_\pi$.
    We show by induction on $\leadsto_\pi$ that when
    $S_A, \Gamma_A \leadsto_\pi S_C, \Gamma_C$ and
    $\exists \Enforcer_A. \forall \sigma \in \Traces. \sigma \in \Language(S_A)$
    then there exist two enforcement event systems
    $\EventSystem_C$ and $\EventSystem_A$ such that
    $\EventSystem_C$ is sound for $S_C$,
    $\EventSystem_A$ is sound for $S_A$, and
    $\EventSystem_C \refineL_{\cConv(\pi)} \EventSystem_A$.

    In the rest of this proof, let
    $R_\pi \subseteq \Traces_{\sys C} \times \Traces_{\sys A}$
    be the simulation relation
    defined by
    $\overrightarrow{\cConv(\pi)}: \Traces_{\sys C} \to \Traces_{\sys A}$.

    \textbf{REFL} In this case $S_C = S_A$ and $\pi$ is the 
    identity mapping.
    By Theorem
    \ref{theo:induced-enforcement-event-system-soundness}
    we know that $\EventSystem(\Enforcer_A)$ is sound for
    $S_A$.
    Finally, any enforcement event system 
    $\EventSystem(\Enforcer_A)$ is trivially a refinement
    of itself under the identity mapping
    (and the identity simulation relation).
    
    \textbf{TRANS} For this case we can simply apply the
    induction hypothesis twice.

    \textbf{REF}
    We know that two enforcers $\Enforcer_C$ that is
    sound for $S_C$ exists
    (Theorems \ref{theo:cre-refinement-enforceability})
    and that its induced event Systems
    are also sound for $S_C$
    (\ref{theo:induced-enforcement-event-system-soundness}).

    It remains to be shown that the two induced
    event systems refine each other.

    Let's consider the situation in point 2 of 
    Definition \ref{def:enforcement-event-system-refinement}.
    We have two concrete system traces
    $\sigma_{\sys C}, \sigma_{\sys C}' \in \Traces_{\sys C}$,
    one abstract system trace
    $\sigma_{\sys A} \in \Traces_{\sys A}$, and a concrete event
    $e_C \in \Eventspace_C$ such that
    $\sigma_{\sys C} \Transitionrelation{C}{e_C} \sigma_{\sys C}'$.
    And if this is the case we want to prove that there
    exists an abstract system trace
    $\sigma_{\sys A}' \in \Traces_{\sys A}$ such that
    $\sigma_{\sys A} \Transitionrelation{A}{\cConv(\pi)(\sigma_{\sys C}, e_C)} \sigma_{\sys A}'$.

    We can argue that the abstract enforcement event system can
    iteratively be defined by applying the mediator function $\cConv(\pi)$
    to all traces and transitions of the concrete enforcement event system.
    A transition 
    $\sigma_{\sys C} \Transitionrelation{C}{e_C} \sigma_{\sys C}'$
    would then become
    $\overrightarrow{\cConv(\pi)}(\sigma_{\sys C})
    \Transitionrelation{A}{\cConv(\pi)(\sigma_{\sys C}, e_C)}
    \overrightarrow{\cConv(\pi)}(\sigma_{\sys C}')$.
    And the simulation relation $R$ can be defined by the aforementioned
    $R_\pi$.
    We call the so constructed enforcement event system $\EventSystem_A$.
    The event system $\EventSystem_A$ is sound for $S_A$ because
    of Theorem \ref{theo:cre-refinement-trace-containment}.

    \textbf{IMP$_R$}
    This rule simply strengthens a policy and thus reduces the
    set of "allowed" traces.
    By Theorem \ref{theo:cre-refinement-enforceability}
    we know that there exists an enforcer $\Enforcer_C$
    that is sound for $S_C$.
    By the fact that $\Language(S_C) \subseteq \Language(S_A)$
    we can deduce that $\Enforcer_C$ is also
    sound for $S_A$.
    And thus the induced event system for $\Enforcer_C$
    is sound for both $S_C$ and $S_A$,
    and also a refinement of itself.

    \textbf{IMP$_C$} and \textbf{IMP$_E$}
    Without loss of generality we will only consider
    the case of IMP$_C$.
    The case of IMP$_E$ is analogous.
    This rule strengthens or weakens the definition of
    an internal event alongside
    the addition of a new, enforceable regulative rule.
    The regulative rule, like in the IMP$_R$ case,
    reduces the set of "allowed" traces.
    The strengthened or weakened definition has
    the effect that the defined event will occur less
    or more often, but under the (anti) monotonicity
    constraints it is ensured that the missing or extra
    occurrences of the event cannot affect compliance.

    The event systems $\EventSystem_C$ and $\EventSystem_A$
    can be simply constructed by inducing them from the
    sound enforcers $\Enforcer_C$ and $\Enforcer_A$
    which we know must exist.
    Then we can build up the simulation relation $R$, by
    iteratively following all possible transitions and adding
    the created trace pairs to $R$.
    e.g. starting with the empty trace, we check what trace
    the concrete enforcement event system
    $\EventSystem_C$ would produce when the event
    $e$ occurs and what trace the abstract enforcement
    event system $\EventSystem_A$ would produce when
    the event $e$ occurs, and then we add the pair of resulting
    traces to $R$.
\end{proof}

% \begin{theorem}
%   \label{thm:refinement_soundness}
%   If $\con{\mathcal{S}} \refines_\pi \abs{\mathcal{S}}$ and $\sigma \vDash \con{\mathcal{S}}$, then $\pi(\sigma) \vDash \abs{\mathcal{S}}$.
% \end{theorem} 

% \begin{theorem}
%   \label{thm:refinement_enforceability}
%   If $\con{\mathcal{S}} \refines_\pi \abs{\mathcal{S}}$ and $\abs{\mathcal{S}}$ is enforceable,
%   then  $\con{\mathcal{S}}$ is enforceable.
% \end{theorem}
